\documentclass[11pt]{article}
\newcommand{\doctitle}{Simplification Plan for pmsimstats-simple}
\newcommand{\shorttitle}{Simplification Plan}
\newcommand{\docdate}{2026-03-24}
% pmsimstats-ng standard document preamble
% Usage: % pmsimstats-ng standard document preamble
% Usage: % pmsimstats-ng standard document preamble
% Usage: \input{pmsimstats-preamble} after \documentclass[11pt]{article}
%
% Documents must define before \input:
%   \newcommand{\doctitle}{...}       Full document title
%   \newcommand{\shorttitle}{...}     Short title for header
%   \newcommand{\docdate}{YYYY-MM-DD} Document date

\usepackage[margin=1in]{geometry}
\usepackage{amsmath,amssymb}
\usepackage[T1]{fontenc}
\usepackage{lmodern}
\usepackage{microtype}
\usepackage{graphicx}
\graphicspath{{figures/}}
\usepackage{booktabs}
\usepackage{longtable}
\usepackage{array}
\usepackage{hyperref}
\usepackage{xcolor}
\usepackage{fancyhdr}
\usepackage{parskip}
\usepackage{caption}
\usepackage{enumitem}
\usepackage{verbatim}
\usepackage{listings}

\lstset{
  language=R,
  basicstyle=\small\ttfamily,
  keywordstyle=\color{blue!70!black},
  commentstyle=\color{green!50!black},
  stringstyle=\color{red!60!black},
  breaklines=true,
  frame=single,
  framerule=0.5pt,
  rulecolor=\color{gray!40},
  backgroundcolor=\color{gray!5},
  xleftmargin=1em,
  framexleftmargin=1em,
  columns=flexible
}

\hypersetup{
  colorlinks=true,
  linkcolor=blue!60!black,
  citecolor=green!50!black,
  urlcolor=blue!60!black
}

\pagestyle{fancy}
\fancyhf{}
\lhead{\footnotesize pmsimstats-ng Technical Report}
\rhead{\footnotesize \shorttitle}
\rfoot{\footnotesize \thepage}
\renewcommand{\headrulewidth}{0.4pt}

\title{\doctitle}
\author{pmsimstats team}
\date{\docdate}
 after \documentclass[11pt]{article}
%
% Documents must define before \input:
%   \newcommand{\doctitle}{...}       Full document title
%   \newcommand{\shorttitle}{...}     Short title for header
%   \newcommand{\docdate}{YYYY-MM-DD} Document date

\usepackage[margin=1in]{geometry}
\usepackage{amsmath,amssymb}
\usepackage[T1]{fontenc}
\usepackage{lmodern}
\usepackage{microtype}
\usepackage{graphicx}
\graphicspath{{figures/}}
\usepackage{booktabs}
\usepackage{longtable}
\usepackage{array}
\usepackage{hyperref}
\usepackage{xcolor}
\usepackage{fancyhdr}
\usepackage{parskip}
\usepackage{caption}
\usepackage{enumitem}
\usepackage{verbatim}
\usepackage{listings}

\lstset{
  language=R,
  basicstyle=\small\ttfamily,
  keywordstyle=\color{blue!70!black},
  commentstyle=\color{green!50!black},
  stringstyle=\color{red!60!black},
  breaklines=true,
  frame=single,
  framerule=0.5pt,
  rulecolor=\color{gray!40},
  backgroundcolor=\color{gray!5},
  xleftmargin=1em,
  framexleftmargin=1em,
  columns=flexible
}

\hypersetup{
  colorlinks=true,
  linkcolor=blue!60!black,
  citecolor=green!50!black,
  urlcolor=blue!60!black
}

\pagestyle{fancy}
\fancyhf{}
\lhead{\footnotesize pmsimstats-ng Technical Report}
\rhead{\footnotesize \shorttitle}
\rfoot{\footnotesize \thepage}
\renewcommand{\headrulewidth}{0.4pt}

\title{\doctitle}
\author{pmsimstats team}
\date{\docdate}
 after \documentclass[11pt]{article}
%
% Documents must define before \input:
%   \newcommand{\doctitle}{...}       Full document title
%   \newcommand{\shorttitle}{...}     Short title for header
%   \newcommand{\docdate}{YYYY-MM-DD} Document date

\usepackage[margin=1in]{geometry}
\usepackage{amsmath,amssymb}
\usepackage[T1]{fontenc}
\usepackage{lmodern}
\usepackage{microtype}
\usepackage{graphicx}
\graphicspath{{figures/}}
\usepackage{booktabs}
\usepackage{longtable}
\usepackage{array}
\usepackage{hyperref}
\usepackage{xcolor}
\usepackage{fancyhdr}
\usepackage{parskip}
\usepackage{caption}
\usepackage{enumitem}
\usepackage{verbatim}
\usepackage{listings}

\lstset{
  language=R,
  basicstyle=\small\ttfamily,
  keywordstyle=\color{blue!70!black},
  commentstyle=\color{green!50!black},
  stringstyle=\color{red!60!black},
  breaklines=true,
  frame=single,
  framerule=0.5pt,
  rulecolor=\color{gray!40},
  backgroundcolor=\color{gray!5},
  xleftmargin=1em,
  framexleftmargin=1em,
  columns=flexible
}

\hypersetup{
  colorlinks=true,
  linkcolor=blue!60!black,
  citecolor=green!50!black,
  urlcolor=blue!60!black
}

\pagestyle{fancy}
\fancyhf{}
\lhead{\footnotesize pmsimstats-ng Technical Report}
\rhead{\footnotesize \shorttitle}
\rfoot{\footnotesize \thepage}
\renewcommand{\headrulewidth}{0.4pt}

\title{\doctitle}
\author{pmsimstats team}
\date{\docdate}


\usepackage{fancyvrb}
\newcolumntype{L}[1]{>{\raggedright\arraybackslash}p{#1}}
\newcolumntype{C}[1]{>{\centering\arraybackslash}p{#1}}

\begin{document}
\maketitle
\thispagestyle{fancy}
\tableofcontents
\newpage

%% ============================================================
\section{Overview}
%% ============================================================

This document outlines possible simplifications to the N-of-1 trial
simulation for biomarker-moderated treatment response. The goal is to
reduce complexity while preserving the core scientific question:
\textbf{What is the statistical power to detect a biomarker $\times$
treatment interaction across different trial designs?}


%% ============================================================
\section{Current Complexity in pmsimstats2025}
%% ============================================================

\begin{table}[h]
\centering
\begin{tabular}{@{}lll@{}}
\toprule
Component & Current Implementation & Complexity \\
\midrule
Covariance structure & Full MVN with AR(1), cross-component correlations & High \\
Data generation & Two-stage conditional sampling from joint distribution & High \\
Analysis model & Mixed-effects (lme) with CAR(1) correlation & High \\
Carryover & Exponential decay via half-life parameter & Medium \\
Censoring & 5 different dropout/censoring patterns & Medium \\
Response components & BR, ER, TR with separate random effects per timepoint & High \\
\bottomrule
\end{tabular}
\end{table}


%% ============================================================
\section{Proposed Simplifications}
%% ============================================================

%% ------------------------------------------------------------
\subsection{Data Generation}
%% ------------------------------------------------------------

\subsubsection{Covariance Structure}

\textbf{Current:} Joint multivariate normal for BR, ER, TR random
effects across all timepoints, plus biomarker and baseline, with AR(1)
temporal structure and cross-component correlations.

\textbf{Simplified options:}

\begin{table}[h]
\centering\small
\begin{tabular}{@{}L{3cm}L{4cm}L{3cm}L{3.5cm}@{}}
\toprule
Option & Description & Pros & Cons \\
\midrule
Independent errors & $\text{response} = \text{signal} + \varepsilon$ & Trivial to implement & Ignores repeated measures \\
Compound symmetry & Equal correlation between all timepoints & Simple, one parameter & Unrealistic for longitudinal \\
Random intercept only & $Y = \text{signal} + u_i + \varepsilon_{ij}$ & Standard mixed model & No temporal correlation \\
AR(1) residuals only & Single AR(1) process, no component separation & Moderate complexity & Loses BR/ER/TR decomposition \\
\bottomrule
\end{tabular}
\end{table}

\textbf{Recommendation:} Start with \textbf{random intercept +
independent errors}. This captures between-subject variability without
temporal complexity.

\begin{Verbatim}[fontsize=\small]
u_i <- rnorm(1, 0, sigma_between)
epsilon <- rnorm(n_timepoints, 0, sigma_within)
response <- baseline + signal + u_i + epsilon
\end{Verbatim}

\subsubsection{Response Components}

\textbf{Current:} Separate BR, ER, TR components each with own random
effects.

\textbf{Simplified options:}

\begin{table}[h]
\centering\small
\begin{tabular}{@{}L{3.5cm}L{4.5cm}L{5.5cm}@{}}
\toprule
Option & Description & Implementation \\
\midrule
Single treatment effect & Collapse BR into one effect & \texttt{effect = on\_drug * beta * (1 + bm\_mod * biomarker)} \\
No ER component & Remove expectancy response & Reduces parameters by 1/3 \\
No TR component & Remove time/regression to mean & Assumes stable baseline \\
Combined ER+TR & Single ``non-drug'' trajectory & \texttt{non\_drug = week * rate} \\
\bottomrule
\end{tabular}
\end{table}

\textbf{Recommendation:} Keep BR with biomarker moderation (core
hypothesis). Consider dropping TR initially; add back if baseline drift
is important.

%% ------------------------------------------------------------
\subsection{Statistical Analysis}
%% ------------------------------------------------------------

\subsubsection{Model Choice}

\textbf{Current:} \texttt{nlme::lme()} with CAR(1) correlation structure.

\textbf{Simplified options:}

\begin{table}[h]
\centering\small
\begin{tabular}{@{}L{2.2cm}L{4.8cm}L{3cm}L{3.5cm}@{}}
\toprule
Option & Formula & Pros & Cons \\
\midrule
ANCOVA & \texttt{lm(post \textasciitilde{} pre + drug * biomarker)} & Fast, familiar, closed-form & Requires summary measure \\
$t$-test on change & \texttt{t.test(delta \textasciitilde{} group)} & Simplest & Ignores biomarker \\
GEE & \texttt{geepack::geeglm()} & Robust to correlation misspec. & Slightly complex \\
Random intercept LMM & \texttt{lme4::lmer(y \textasciitilde{} x + (1|id))} & Simpler than full lme & Still mixed model \\
Simple linear regression & \texttt{lm(response \textasciitilde{} on\_drug * biomarker)} & Ignores clustering & Biased SE but fast \\
RM-ANOVA & \texttt{aov(y \textasciitilde{} bm\_group * drug + Error(id/drug))} & Familiar to clinicians & Requires dichotomized biomarker ($\sim$36\% information loss); sphericity assumption \\
\bottomrule
\end{tabular}
\end{table}

\textbf{Recommendation:} Use \textbf{ANCOVA on change scores} as
primary analysis:

\begin{Verbatim}[fontsize=\small]
summary_data <- trial_data |>
  group_by(participant_id, on_drug) |>
  summarise(mean_response = mean(response), .groups = "drop") |>
  pivot_wider(names_from = on_drug, values_from = mean_response) |>
  mutate(change = `TRUE` - `FALSE`)

model <- lm(change ~ biomarker, data = summary_data)
p_value <- coef(summary(model))["biomarker", "Pr(>|t|)"]
\end{Verbatim}

\subsubsection{Handling Repeated Measures}

\textbf{Current:} Full longitudinal model with all timepoints.

\textbf{Simplified options:}

\begin{table}[h]
\centering
\begin{tabular}{@{}ll@{}}
\toprule
Option & Description \\
\midrule
Summary statistics & Mean/AUC during each phase \\
Last observation & Use final timepoint only \\
First vs.\ last & Compare endpoints \\
Phase means & Average within treatment phases \\
\bottomrule
\end{tabular}
\end{table}

\textbf{Recommendation:} \textbf{Phase means} -- average response during
on-drug vs.\ off-drug periods, then analyze the difference.


%% ============================================================
\subsection{Mathematical Specification of the Interaction Test}
\label{sec:mathspec}
%% ============================================================

This section formally defines the biomarker-treatment interaction test
under each analysis method and connects the precision of the interaction
estimate to trial design features: the number of subjects ($N$), the
number of assessment timepoints per subject ($T$), the number of
within-subject treatment transitions (crossovers), and the proportion of
time spent on vs.\ off drug.

\subsubsection{Notation}

\begin{table}[h]
\centering\small
\begin{tabular}{@{}cl@{}}
\toprule
Symbol & Definition \\
\midrule
$Y_{it}$ & Observed outcome for subject $i$ at timepoint $t$ \\
$bm_i$ & Baseline biomarker value for subject $i$ (time-invariant) \\
$D_{it}$ & Drug exposure indicator ($1 = $ on drug, $0 = $ off drug) \\
$t$ & Time (weeks since baseline) \\
$u_i$ & Subject-specific random intercept, $u_i \sim N(0, \sigma^2_u)$ \\
$\varepsilon_{it}$ & Residual error, $\varepsilon_{it} \sim N(0, \sigma^2_\varepsilon)$ \\
$N$ & Number of subjects \\
$T_i$ & Total assessment timepoints for subject $i$ \\
$n_{i,\text{on}}$ & Number of on-drug assessments for subject $i$ \\
$n_{i,\text{off}}$ & Number of off-drug assessments for subject $i$ \\
$K_i$ & Number of treatment transitions for subject $i$ \\
\bottomrule
\end{tabular}
\end{table}

Throughout, the \textbf{target hypothesis} is:
\[
H_0: \text{the biomarker does not predict the magnitude of the
drug-specific response.}
\]
The mathematical form of this hypothesis differs depending on whether
the design includes within-subject drug variation.

\subsubsection{The Fundamental Distinction: Within-Subject Drug Variation}

\textbf{Designs with within-subject drug variation} (OL+BDC, CO,
N-of-1): Subject $i$ has observations both on drug ($D_{it} = 1$) and
off drug ($D_{it} = 0$). The interaction is $bm_i \times D_{it}$,
testing whether the biomarker predicts the within-subject drug effect.

\textbf{Designs without within-subject drug variation} (OL): All
observations for subject $i$ have $D_{it} = 1$. There is no
within-subject drug contrast. The interaction must be reformulated as
$bm_i \times t$, testing whether the biomarker predicts the rate of
improvement over time while on drug. This is a fundamentally weaker test
because temporal improvement confounds BR, ER, and TV.

This distinction is the primary reason the open-label design has low
power (Hendrickson Figure~4A): the interaction term targets a different,
noisier quantity.

%% --- Method 1 ---
\subsubsection{Method 1: ANCOVA on Phase-Mean Change Scores}

\textbf{Step 1.} Collapse repeated measures into within-subject phase
means:
\[
\bar{Y}_{i,\text{on}} = \frac{1}{n_{i,\text{on}}}
\sum_{t:\, D_{it}=1} Y_{it}, \qquad
\bar{Y}_{i,\text{off}} = \frac{1}{n_{i,\text{off}}}
\sum_{t:\, D_{it}=0} Y_{it}
\]

\textbf{Step 2.} Compute the within-subject drug effect:
\[
\Delta_i = \bar{Y}_{i,\text{on}} - \bar{Y}_{i,\text{off}}
\]

\textbf{Step 3.} Regress the drug effect on the biomarker:
\[
\Delta_i = \gamma_0 + \gamma_1 \cdot bm_i + e_i
\]

\textbf{Interaction test:} $H_0: \gamma_1 = 0$ via $t$-test on
$\hat{\gamma}_1$, with $N - 2$ degrees of freedom.

\textbf{Precision of $\hat{\gamma}_1$:}
\[
\operatorname{Var}(\hat{\gamma}_1) =
\frac{\sigma^2_\Delta}{N \cdot \operatorname{Var}(bm)}
\]
where $\sigma^2_\Delta = \operatorname{Var}(\Delta_i)$ is the
between-subject variance of the within-subject drug effect. This
variance decomposes as:
\[
\sigma^2_\Delta = \operatorname{Var}(\bar{Y}_{i,\text{on}}) +
\operatorname{Var}(\bar{Y}_{i,\text{off}}) -
2\operatorname{Cov}(\bar{Y}_{i,\text{on}},\bar{Y}_{i,\text{off}})
\]
Each phase-mean variance is approximately:
\[
\operatorname{Var}(\bar{Y}_{i,\text{phase}}) \approx
\sigma^2_u + \frac{\sigma^2_\varepsilon}{n_{i,\text{phase}}}
\]
The covariance term arises from the shared random intercept $u_i$ and
equals $\sigma^2_u$ under compound symmetry. Substituting:
\[
\sigma^2_\Delta \approx
\sigma^2_\varepsilon \!\left(
\frac{1}{n_{i,\text{on}}} + \frac{1}{n_{i,\text{off}}}
\right)
\]
The random intercept cancels in the difference, which is a key advantage
of within-subject contrasts. What remains is purely residual variance,
reduced by the number of observations in each phase.

\textbf{Design implications:}

\begin{table}[h]
\centering\small
\begin{tabular}{@{}L{4.5cm}L{9cm}@{}}
\toprule
Design feature & Effect on $\operatorname{Var}(\hat{\gamma}_1)$ \\
\midrule
More subjects ($N \uparrow$) & Direct reduction: $\propto 1/N$ \\
More on-drug assessments ($n_{\text{on}} \uparrow$) & Reduces $\sigma^2_\Delta$ via $1/n_{\text{on}}$ term \\
More off-drug assessments ($n_{\text{off}} \uparrow$) & Reduces $\sigma^2_\Delta$ via $1/n_{\text{off}}$ term \\
Balanced phases ($n_{\text{on}} \approx n_{\text{off}}$) & Minimizes the harmonic mean, hence $\sigma^2_\Delta$ \\
More crossovers ($K \uparrow$) & No direct effect beyond increasing $n_{\text{on}}$ and $n_{\text{off}}$; distributing observations across more transitions reduces sensitivity to period effects \\
Wider biomarker range ($\operatorname{Var}(bm) \uparrow$) & Direct reduction: $\propto 1/\operatorname{Var}(bm)$ \\
\bottomrule
\end{tabular}
\end{table}

\textbf{Design-specific behavior:}
\begin{itemize}[nosep]
\item \textbf{OL:} $n_{i,\text{off}} = 0$, so $\Delta_i$ is undefined.
  The ANCOVA must be reformulated as
  $\bar{Y}_i = \gamma_0 + \gamma_1 \cdot bm_i + e_i$ (regress total
  response on biomarker), which confounds BR with ER and TV.
\item \textbf{OL+BDC:} $n_{i,\text{on}} \approx 5$,
  $n_{i,\text{off}} \approx 2$. Imbalanced phases inflate
  $\sigma^2_\Delta$.
\item \textbf{CO:} $n_{i,\text{on}} \approx 4$,
  $n_{i,\text{off}} \approx 4$. Balanced phases. One crossover
  transition.
\item \textbf{N-of-1:} $n_{i,\text{on}} \approx 6$,
  $n_{i,\text{off}} \approx 4$. Multiple crossover transitions. Highest
  total observation count and reasonable balance.
\end{itemize}

\textbf{OL special case.} Without drug variation, the ANCOVA tests a
different quantity. One approach is to use the slope of improvement over
time as the subject-level summary:
\[
b_i = \text{slope of } Y_{it} \text{ vs.\ } t \text{ for subject } i
\]
\[
b_i = \gamma_0 + \gamma_1 \cdot bm_i + e_i
\]
Here $\gamma_1$ tests whether the biomarker predicts the rate of
improvement. This conflates BR rate with ER and TV rates, explaining the
poor power of the OL design.

%% --- Method 2 ---
\subsubsection{Method 2: Paired $t$-test on Change Scores}

The paired $t$-test applies to the same $\Delta_i$ as ANCOVA but tests
a different hypothesis:
\[
H_0: E[\Delta_i] = 0 \quad \text{(no average drug effect)}
\]
\[
t = \frac{\bar{\Delta}}{s_\Delta / \sqrt{N}}, \quad df = N - 1
\]

\textbf{This does not test the biomarker interaction.} It tests whether
the drug works on average, collapsing over biomarker values. To recover
a biomarker interaction test from a $t$-test framework, one must
dichotomize:
\[
\text{Split subjects into } bm_{\text{high}} \text{ and }
bm_{\text{low}} \text{ by median}
\]
\[
t = \frac{\bar{\Delta}_{\text{high}} -
\bar{\Delta}_{\text{low}}}{s_p \sqrt{1/N_{\text{high}} +
1/N_{\text{low}}}}, \quad df = N - 2
\]

This is equivalent to ANCOVA with a binary biomarker and inherits the
same design dependencies. Dichotomization discards information, reducing
power relative to ANCOVA with a continuous biomarker by approximately
$2/\pi \approx 64\%$ relative efficiency (Cohen 1983).

\textbf{Design implications:} Same as ANCOVA, but uniformly less
powerful due to dichotomization. Not recommended when the biomarker is
continuous.

%% --- Method 3 ---
\subsubsection{Method 3: GEE (Generalized Estimating Equations)}

\textbf{Model (population-averaged):}
\[
E[Y_{it}] = \beta_0 + \beta_1 \cdot bm_i + \beta_2 \cdot D_{it}
+ \beta_3 \cdot t + \beta_4 \cdot bm_i \cdot D_{it}
\]
with working correlation matrix $R_i(\alpha)$ and sandwich (robust)
variance estimator.

\textbf{Interaction test:} $H_0: \beta_4 = 0$ via Wald test using the
sandwich-estimated variance:
\[
z = \frac{\hat{\beta}_4}{\widehat{\text{SE}}_{\text{robust}}
(\hat{\beta}_4)}
\]

\textbf{Precision of $\hat{\beta}_4$:}

The sandwich estimator is consistent regardless of the working
correlation structure, but efficiency depends on how well the working
correlation approximates the truth. Under a correct working correlation:
\[
\operatorname{Var}(\hat{\beta}_4) \approx
\frac{\sigma^2_{\text{eff}}}{N \cdot \operatorname{Var}(bm) \cdot
\operatorname{Var}_{\text{within}}(D)}
\]
where $\operatorname{Var}_{\text{within}}(D)$ is the within-subject
variance of drug exposure and $\sigma^2_{\text{eff}}$ is the effective
residual variance after accounting for the working correlation.

The \textbf{within-subject variance of drug exposure} is the critical
design-dependent quantity:
\[
\operatorname{Var}_{\text{within}}(D_i) = \bar{D}_i(1 - \bar{D}_i)
\]
where $\bar{D}_i = n_{i,\text{on}} / T_i$ is the proportion of time on
drug. This is maximized when $\bar{D}_i = 0.5$ (equal time on and off
drug) and equals zero when $\bar{D}_i = 1$ (the OL design).

\textbf{Design implications:}

\begin{table}[h]
\centering\small
\begin{tabular}{@{}L{5cm}L{8.5cm}@{}}
\toprule
Design feature & Effect on $\operatorname{Var}(\hat{\beta}_4)$ \\
\midrule
More subjects ($N \uparrow$) & Direct reduction: $\propto 1/N$ \\
More timepoints ($T \uparrow$) & Reduces $\sigma^2_{\text{eff}}$; modest gain under strong autocorrelation \\
Balanced drug exposure ($\bar{D} \to 0.5$) & Maximizes $\operatorname{Var}_{\text{within}}(D)$, directly reducing $\operatorname{Var}(\hat{\beta}_4)$ \\
More crossovers ($K \uparrow$) & Does not change $\operatorname{Var}_{\text{within}}(D)$ directly, but distributes drug variation across time \\
Autocorrelation ($\rho \uparrow$) & Increases effective variance; nearby observations contribute less independent information \\
\bottomrule
\end{tabular}
\end{table}

\textbf{Design-specific behavior:}
\begin{itemize}[nosep]
\item \textbf{OL:} $\operatorname{Var}_{\text{within}}(D) = 0$. The
  interaction $bm \times D$ is not estimable. Must substitute
  $bm \times t$.
\item \textbf{OL+BDC:} $\bar{D} \approx 0.75$.
  $\operatorname{Var}_{\text{within}}(D) \approx 0.19$. Suboptimal but
  nonzero.
\item \textbf{CO:} $\bar{D} = 0.5$.
  $\operatorname{Var}_{\text{within}}(D) = 0.25$. Optimal balance.
\item \textbf{N-of-1:} $\bar{D} \approx 0.55$.
  $\operatorname{Var}_{\text{within}}(D) \approx 0.25$. Near-optimal,
  with multiple transitions.
\end{itemize}

\textbf{GEE vs.\ LMM:} GEE estimates population-averaged effects
(marginal model), while LMM estimates subject-specific effects
(conditional model). For the linear case with continuous outcomes,
$\beta_4$ has the same interpretation under both approaches. The
practical difference is that GEE provides valid inference under
working-correlation misspecification (via the sandwich estimator), while
LMM requires correct specification of the random effects structure but
is more efficient when correctly specified.

%% --- Method 4 ---
\subsubsection{Method 4: Random Intercept LMM (\texttt{lme4::lmer})}

\textbf{Model (subject-specific):}
\[
Y_{it} = \beta_0 + \beta_1 \cdot bm_i + \beta_2 \cdot D_{it}
+ \beta_3 \cdot t + \beta_4 \cdot bm_i \cdot D_{it}
+ u_i + \varepsilon_{it}
\]
\[
u_i \sim N(0, \sigma^2_u), \qquad
\varepsilon_{it} \sim N(0, \sigma^2_\varepsilon)
\]

\textbf{Interaction test:} $H_0: \beta_4 = 0$ via Wald test (or
Kenward-Roger / Satterthwaite $F$-test for small $N$).

\textbf{Precision of $\hat{\beta}_4$:}

Under the random-intercept model with independent residuals, the
information matrix for $\beta_4$ depends on the effective within-subject
information after projecting out the random intercept:
\[
\operatorname{Var}(\hat{\beta}_4) \approx
\frac{\sigma^2_\varepsilon}{
\displaystyle\sum_{i=1}^{N} \widetilde{bm}_i^{\,2} \cdot
\text{SS}_{\text{within}}(D_i)}
\]
where $\widetilde{bm}_i = bm_i - \overline{bm}$ is the centered
biomarker value and:
\[
\text{SS}_{\text{within}}(D_i) =
\sum_{t=1}^{T_i} (D_{it} - \bar{D}_i)^2
= T_i \cdot \bar{D}_i (1 - \bar{D}_i)
\]
is the within-subject sum of squares of drug exposure for subject $i$.
This quantity is the total ``information'' about drug variation contained
in subject $i$'s data.

\textbf{Design implications:}

\begin{table}[h]
\centering\small
\begin{tabular}{@{}L{3.5cm}L{4.5cm}L{5.5cm}@{}}
\toprule
Design feature & Effect on $\operatorname{Var}(\hat{\beta}_4)$ & Mechanism \\
\midrule
$N \uparrow$ & $\propto 1/N$ & More independent contrasts \\
$T \uparrow$ & $\propto 1/T$ (given fixed $\bar{D}$) & Each timepoint adds to $\text{SS}_{\text{within}}(D_i)$ \\
$\bar{D} \to 0.5$ & Maximizes $\bar{D}(1-\bar{D})$ & Equal on/off time maximizes contrast \\
$K \uparrow$ & No direct effect under independence & Under AR(1), more transitions reduce autocorrelation penalty \\
Wider biomarker range & Increases $\sum \widetilde{bm}_i^{\,2}$ & More leverage for the interaction \\
$\rho > 0$ & Inflates effective $\sigma^2_\varepsilon$ & Correlated residuals reduce effective sample size \\
\bottomrule
\end{tabular}
\end{table}

\textbf{The autocorrelation-crossover interaction.} Under AR(1)
residuals (as implemented in the full \texttt{nlme::lme} model with
\texttt{corCAR1}), the effective information from each on-drug or
off-drug block depends on block length $L$ and autocorrelation $\rho$:
\[
\text{effective observations per block} \approx
\frac{L}{1 + (L-1)\rho} \cdot L
\]
Short blocks (high $K$, low $L$) suffer less from autocorrelation than
long blocks (low $K$, high $L$) because fewer consecutive observations
are highly correlated. This is why the N-of-1 design (four 4-week
blocks) can outperform the crossover (two 10-week blocks) when
autocorrelation is moderate.

Conversely, when the drug has substantial carryover, each transition
introduces a contaminated observation at the block boundary. More
transitions ($K \uparrow$) means more contaminated observations. This
is the mechanism behind the ``precipitous decline in power'' with
carryover that Hendrickson et al.\ report (their Figure~4B).

\textbf{Design-specific $\text{SS}_{\text{within}}(D)$:}

\begin{table}[h]
\centering
\begin{tabular}{@{}lcccccc@{}}
\toprule
Design & $T$ & $n_{\text{on}}$ & $n_{\text{off}}$ & $\bar{D}$ &
$\text{SS}_{\text{within}}(D)$ & $K$ \\
\midrule
OL & 8 & 8 & 0 & 1.0 & 0 & 0 \\
OL+BDC & 8 & ${\sim}6$ & ${\sim}2$ & ${\sim}0.75$ & ${\sim}1.5$ & 1 \\
CO & 8 & 4 & 4 & 0.5 & 2.0 & 1 \\
N-of-1 & 12 & ${\sim}7$ & ${\sim}5$ & ${\sim}0.58$ & ${\sim}2.9$ & 3 \\
\bottomrule
\end{tabular}
\end{table}

The N-of-1 design has the highest $\text{SS}_{\text{within}}(D)$ due to
both more total timepoints and near-balanced drug exposure. The OL
design has zero within-subject drug information.

%% --- Method 5 ---
\subsubsection{Method 5: Simple Linear Regression (Ignoring Clustering)}

\textbf{Model:}
\[
Y_{it} = \beta_0 + \beta_1 \cdot bm_i + \beta_2 \cdot D_{it}
+ \beta_3 \cdot t + \beta_4 \cdot bm_i \cdot D_{it}
+ \varepsilon_{it}
\]
This is the same fixed-effects structure as Method~4 but without the
random intercept $u_i$. All $N \times T$ observations are treated as
independent.

\textbf{Interaction test:} $H_0: \beta_4 = 0$ via OLS $t$-test.

\textbf{Precision of $\hat{\beta}_4$ (nominal):}
\[
\operatorname{Var}_{\text{OLS}}(\hat{\beta}_4) =
\frac{\hat{\sigma}^2}{\text{SS}(bm \cdot D \mid
\text{other predictors})}
\]
where $\hat{\sigma}^2$ is the OLS residual variance. Because the model
ignores the random intercept, $\hat{\sigma}^2$ absorbs both
$\sigma^2_u$ and $\sigma^2_\varepsilon$, but the denominator treats all
$N \times T$ observations as independent. This produces
\textbf{downward-biased standard errors} because the effective sample
size for the between-subject component ($bm_i$) is $N$, not $N \times T$.

\textbf{Type I error inflation:} The degree of inflation depends on
$\sigma^2_u / \sigma^2_\varepsilon$ (the intraclass correlation) and
$T$. With $T = 8$ and ICC $= 0.3$, the nominal $\alpha = 0.05$ test
has an actual Type~I error of approximately 0.15--0.25.

\textbf{Design implications:} Same as Method~4 for power ordering
across designs, but with inflated Type~I error that invalidates nominal
$p$-values. Not recommended for inference but can serve as a fast
screening tool.

%% --- Method 6 ---
\subsubsection{Method 6: Repeated Measures ANOVA (Biomarker Dichotomized)}

\textbf{Setup.} Because repeated measures ANOVA requires a grouping
factor (not a continuous covariate), the biomarker must be dichotomized:
\[
G_i = \begin{cases} 1 & \text{if } bm_i > \text{median}(bm) \\
0 & \text{otherwise} \end{cases}
\]
The data are arranged as a $N \times T$ matrix of outcomes, with two
between-subject groups ($G = 0, 1$) and a within-subject factor $D$
(on-drug vs.\ off-drug). If the design has multiple timepoints within
each drug condition, there is an additional within-subject factor for
time nested within drug phase.

\textbf{Model (split-plot / mixed ANOVA):}
\[
Y_{it} = \mu + \alpha_{G_i} + \delta_{D_{it}} +
(\alpha\delta)_{G_i, D_{it}} + \pi_i + \varepsilon_{it}
\]
where $\mu$ is the grand mean, $\alpha_{G_i}$ is the between-subject
biomarker group effect, $\delta_{D_{it}}$ is the within-subject drug
effect, $(\alpha\delta)_{G_i, D_{it}}$ is the biomarker-by-drug
interaction, $\pi_i \sim N(0, \sigma^2_\pi)$ is the subject effect
(random), and $\varepsilon_{it} \sim N(0, \sigma^2_\varepsilon)$ is the
residual.

\textbf{Interaction test:} $H_0: (\alpha\delta)_{G,D} = 0$ for all
combinations, tested via:
\[
F = \frac{\text{MS}_{G \times D}}{\text{MS}_{\text{error(within)}}},
\quad df_1 = (g-1)(d-1), \quad df_2 = (N-g)(d-1)
\]
where $g = 2$ (biomarker groups) and $d = 2$ (on-drug, off-drug), giving
$df_1 = 1$ and $df_2 = N - 2$.

\textbf{Sphericity.} When more than two levels of the within-subject
factor are present (e.g., multiple timepoints within each drug phase),
the $F$-test assumes sphericity: equal variances of all pairwise
differences between within-subject levels. Violation inflates the
Type~I error rate. The Greenhouse-Geisser (GG) and Huynh-Feldt (HF)
corrections multiply both $df_1$ and $df_2$ by a correction factor
$\hat{\varepsilon} \in [1/(k-1),\, 1]$:
\[
df_1' = df_1 \cdot \hat{\varepsilon}, \qquad
df_2' = df_2 \cdot \hat{\varepsilon}
\]
When the within-subject factor is simply on-drug vs.\ off-drug
($k = 2$), sphericity is satisfied trivially (there is only one
pairwise difference), and no correction is needed.

\textbf{Precision of the interaction test.} For the simple case
($g = 2$ groups, $d = 2$ drug conditions, collapsing timepoints within
phases into means):
\[
\text{MS}_{G \times D} =
\frac{N}{4} \left(
\bar{\Delta}_{\text{high}} - \bar{\Delta}_{\text{low}}
\right)^2
\]
\[
\text{MS}_{\text{error(within)}} =
\frac{1}{N-2} \sum_{i=1}^{N}
(\Delta_i - \bar{\Delta}_{G_i})^2
\]
where $\Delta_i = \bar{Y}_{i,\text{on}} - \bar{Y}_{i,\text{off}}$ as
in the ANCOVA approach. The noncentrality parameter for the $F$-test is:
\[
\lambda = \frac{N}{4} \cdot
\frac{(\mu_{\Delta,\text{high}} -
\mu_{\Delta,\text{low}})^2}{\sigma^2_\Delta}
\]

\textbf{Comparison to ANCOVA (Method~1).} The repeated measures ANOVA
on phase means with dichotomized biomarker is algebraically equivalent
to a two-sample $t$-test on $\Delta_i$ between biomarker groups. The
ANCOVA on change scores with continuous biomarker tests the same
conceptual hypothesis but retains the full biomarker distribution. The
efficiency loss from dichotomization is well established: for a normally
distributed biomarker, the median split retains only
$2/\pi \approx 63.7\%$ of the information, corresponding to a relative
efficiency of approximately 0.64 (Cohen 1983, Royston et al.\ 2006). In
sample-size terms, the repeated measures ANOVA requires approximately
$N/0.64 \approx 1.57N$ subjects to achieve the same power as ANCOVA
with a continuous biomarker.

\textbf{When timepoints are not collapsed.} If the full set of $T$
timepoints is retained as within-subject levels, the model becomes a
two-way mixed ANOVA with $G$ (between) and timepoint (within). The
interaction $G \times \text{timepoint}$ tests whether the two biomarker
groups have different temporal profiles. This formulation:
\begin{itemize}[nosep]
\item Gains power from using all observations (no information loss from
  averaging)
\item Loses specificity by testing a composite hypothesis across all
  timepoints rather than the drug contrast specifically
\item Requires the sphericity assumption across all $T$ timepoints,
  which is typically violated in longitudinal data and necessitates
  GG/HF correction
\item The corrected degrees of freedom can substantially reduce power
  when $\hat{\varepsilon}$ is small (strong autocorrelation, many
  timepoints)
\end{itemize}

\textbf{Design implications:}

\begin{table}[h]
\centering\small
\begin{tabular}{@{}L{4cm}L{4.5cm}L{5cm}@{}}
\toprule
Design feature & Effect on power & Mechanism \\
\midrule
$N \uparrow$ & $\propto N$ (via $\lambda$) & More independent group comparisons \\
More on/off-drug assessments & Reduces $\sigma^2_\Delta$ when collapsing to phase means & More precise $\Delta_i$ \\
Balanced phases & Same benefit as ANCOVA & Minimizes $\operatorname{Var}(\Delta_i)$ \\
Balanced biomarker groups & Maximizes $F$-test power & Median split guarantees this \\
$K \uparrow$ & Indirect benefit via $n_{\text{on}}, n_{\text{off}}$ & Same as ANCOVA \\
Timepoints retained as levels & Increases $df_1$ but triggers sphericity issues & GG/HF correction can offset gains \\
$\rho \uparrow$ & Reduces $\hat{\varepsilon}$ when timepoints retained & Corrected $df$ shrink \\
\bottomrule
\end{tabular}
\end{table}

\textbf{Design-specific behavior:}
\begin{itemize}[nosep]
\item \textbf{OL:} No within-subject drug factor. The within-subject
  factor becomes time, and the interaction $G \times \text{time}$ tests
  whether biomarker groups have different slopes. Same confounding
  problem as all other methods.
\item \textbf{OL+BDC:} The within-subject drug factor has two levels
  but is heavily imbalanced. The sparse off-drug phase inflates
  $\sigma^2_\Delta$.
\item \textbf{CO:} Clean two-level within-subject drug factor with
  balanced phases. Sphericity trivially satisfied. The RM-ANOVA is most
  natural here.
\item \textbf{N-of-1:} Multiple drug-on and drug-off blocks. If
  collapsed to phase means, behaves like ANCOVA with good balance. If
  blocks are retained as separate levels, sphericity becomes a concern.
\end{itemize}

\textbf{Summary.} Repeated measures ANOVA is conceptually accessible and
widely implemented, making it attractive for clinical audiences. However,
for the biomarker interaction test specifically, it suffers from two
compounding limitations: (1)~dichotomization of the biomarker discards
approximately 36\% of the information, and (2)~when timepoints are
retained rather than collapsed, sphericity violations and GG/HF
corrections can further erode power. For these reasons, ANCOVA on change
scores (Method~1) or the random intercept LMM (Method~4) are generally
preferable when the biomarker is continuous.

%% --- OL special case ---
\subsubsection{The OL Design: A Fundamentally Different Interaction Test}

For the open-label design, all six methods face the same problem: there
is no within-subject drug variation. The interaction must be formulated
as biomarker-by-time:
\[
Y_{it} = \beta_0 + \beta_1 \cdot bm_i + \beta_2 \cdot t
+ \beta_3 \cdot bm_i \cdot t + u_i + \varepsilon_{it}
\]
Here $\beta_3$ tests whether the biomarker predicts the slope of
improvement over time. This is a weaker test than $\beta_4$ (the
biomarker-by-drug interaction) for two reasons:
\begin{enumerate}[nosep]
\item \textbf{Confounding.} The slope $dY/dt$ includes BR, ER, and TV.
  The biomarker-by-time interaction conflates all three.
\item \textbf{Precision.} The within-subject information for the time
  slope depends on
  $\text{SS}_{\text{within}}(t) = \sum(t - \bar{t})^2$, which is
  typically smaller than $\text{SS}_{\text{within}}(D)$ for designs
  with balanced crossovers.
\end{enumerate}

%% --- Summary table ---
\subsubsection{Summary: Interaction Test Precision Across Designs and Methods}

Table~\ref{tab:summary} summarizes the effective information for the
biomarker interaction test across all design-method combinations. Entries
are proportional to $1/\operatorname{Var}(\hat{\beta}_{\text{interaction}})$
(higher = more power).

\begin{table}[h]
\centering\small
\caption{Interaction test properties by design and method.}
\label{tab:summary}
\begin{tabular}{@{}L{3.2cm}L{2.8cm}L{2.8cm}L{2.1cm}L{2.5cm}@{}}
\toprule
 & OL & OL+BDC & CO & N-of-1 \\
\midrule
\textbf{Interaction term} & $bm \times t$ & $bm \times D$ & $bm \times D$ & $bm \times D$ \\
\textbf{Target} & Total improvement rate & Drug-specific effect & Drug-specific effect & Drug-specific effect \\
\textbf{Confounded with} & ER, TV slopes & -- & -- & -- \\
\textbf{ANCOVA info} & $N \!\cdot\! \operatorname{Var}(bm) \!\cdot\! \operatorname{Var}(b_i)^{-1}$ & $N \!\cdot\! \operatorname{Var}(bm) \!\cdot\! \sigma^{-2}_\Delta$ & $N \!\cdot\! \operatorname{Var}(bm) \!\cdot\! \sigma^{-2}_\Delta$ & $N \!\cdot\! \operatorname{Var}(bm) \!\cdot\! \sigma^{-2}_\Delta$ \\
\textbf{LMM info} & $\sum \widetilde{bm}^2_i \!\cdot\! \text{SS}(t_i)$ & $\sum \widetilde{bm}^2_i \!\cdot\! 1.5$ & $\sum \widetilde{bm}^2_i \!\cdot\! 2.0$ & $\sum \widetilde{bm}^2_i \!\cdot\! 2.9$ \\
\textbf{GEE info} & $\approx$ LMM & $\approx$ LMM & $\approx$ LMM & $\approx$ LMM \\
\textbf{RM-ANOVA info} & ${\approx}\,0.64 \times$ ANCOVA & ${\approx}\,0.64 \times$ ANCOVA & ${\approx}\,0.64 \times$ ANCOVA & ${\approx}\,0.64 \times$ ANCOVA \\
\textbf{$K$ (transitions)} & 0 & 1 & 1 & 3 \\
\textbf{Autocorr.\ vuln.} & Low & Medium & Medium & Low per block, more boundary effects \\
\textbf{Carryover vuln.} & None & Low (1 transition) & Low (1 transition) & High (3 transitions) \\
\bottomrule
\end{tabular}
\end{table}

The RM-ANOVA row reflects the ${\approx}\,2/\pi$ efficiency penalty from
dichotomizing the biomarker at the median. If timepoints within phases
are retained as separate within-subject levels rather than collapsed to
phase means, GG/HF corrections further reduce effective degrees of
freedom under autocorrelation.

The N-of-1 design achieves the highest information under the LMM and GEE
when carryover is absent, owing to both its larger total $T$ and
near-balanced drug exposure. The crossover achieves optimal balance
($\bar{D} = 0.5$) with fewer transitions, making it more robust to
carryover at the cost of fewer total timepoints.


%% ============================================================
\subsection{Trial Designs (Non-Simplifiable)}
%% ============================================================

\textbf{Critical:} The four Hendrickson trial designs must be preserved
exactly as specified. The design structure is the core scientific
comparison and cannot be simplified.

\begin{table}[h]
\centering\small
\begin{tabular}{@{}llll@{}}
\toprule
Design & Weeks & Structure & Blinding \\
\midrule
Open-Label (OL) & 1--8 & All on drug, unblinded & None \\
OL + Blinded Disc.\ (OL+BDC) & 1--4 OL, 5--8 randomized & Partial \\
Traditional Crossover (CO) & 4 wks drug, 4 wks placebo (order rand.) & Full \\
N-of-1 (Hybrid) & 1--4 OL, 5--12 two 4-wk crossover cycles & Partial \\
\bottomrule
\end{tabular}
\end{table}

Each design must specify: (1)~week-by-week \texttt{on\_drug} indicator,
(2)~blinding status (affects expectancy response), and
(3)~randomization path (which weeks are randomized).

\begin{Verbatim}[fontsize=\small]
designs <- list(
  ol = list(weeks = 1:8,
    on_drug = rep(TRUE, 8), blinded = rep(FALSE, 8),
    expectancy = rep(TRUE, 8)),
  ol_bdc = list(weeks = 1:8,
    on_drug = c(rep(TRUE, 4), NA, NA, NA, NA),
    blinded = c(rep(FALSE, 4), rep(TRUE, 4)),
    expectancy = c(rep(TRUE, 4), rep(FALSE, 4))),
  crossover = list(weeks = 1:8,
    on_drug = NA, blinded = rep(TRUE, 8),
    expectancy = rep(FALSE, 8)),
  nof1 = list(weeks = 1:12,
    on_drug = c(rep(TRUE, 4), rep(NA, 8)),
    blinded = c(rep(FALSE, 4), rep(TRUE, 8)),
    expectancy = c(rep(TRUE, 4), rep(FALSE, 8))))
\end{Verbatim}

\texttt{NA} values indicate randomized assignment at simulation time.


%% ============================================================
\subsection{Carryover Effects}
%% ============================================================

\textbf{Current:} Exponential decay with half-life parameter.

\textbf{Simplified options:}

\begin{table}[h]
\centering\small
\begin{tabular}{@{}L{3cm}L{5.5cm}L{4.5cm}@{}}
\toprule
Option & Formula & Description \\
\midrule
No carryover & $\text{carryover} = 0$ & Assume instant washout \\
Binary carryover & $\mathbf{1}[t_{\text{since}} < \text{washout}]$ & All-or-nothing \\
Linear decay & $\max(0, 1 - t_{\text{since}}/\text{washout})$ & Linear decrease \\
Step function & $0.5 \cdot \mathbf{1}[t_{\text{since}} < 1]$ & Partial for one period \\
\bottomrule
\end{tabular}
\end{table}

\textbf{Recommendation:} Start with \textbf{no carryover}, add
\textbf{linear decay} as sensitivity analysis.


%% ============================================================
\subsection{Censoring/Dropout}
%% ============================================================

\textbf{Current:} Five patterns (none, balanced, high dropout, more
biased, more flat) with complex probability functions.

\textbf{Simplified options:}

\begin{table}[h]
\centering
\begin{tabular}{@{}ll@{}}
\toprule
Option & Description \\
\midrule
No dropout & Complete data for all \\
MCAR & Random dropout, constant probability \\
Single MAR & One dropout pattern tied to response \\
\bottomrule
\end{tabular}
\end{table}

\textbf{Recommendation:} Start with \textbf{no dropout}. Add MCAR
(e.g., 10\% per timepoint) as sensitivity.


%% ============================================================
\subsection{Parameter Space}
%% ============================================================

\begin{table}[h]
\centering\small
\begin{tabular}{@{}lll@{}}
\toprule
Parameter & Current & Simplified \\
\midrule
BR\_rate & 0.05, 0.15, 0.3 & 0.1, 0.3 (2 levels) \\
ER\_rate & 0.05, 0.15, 0.3 & 0.1 (fixed) \\
TR\_rate & 0.05, 0.15, 0.3 & 0 (removed) \\
biomarker\_mod & 0, 0.3, 0.6 & 0.3 (fixed), vary for sensitivity \\
carryover & 0, 0.1, 0.2 & 0 (default), 0.2 (sensitivity) \\
censoring & 5 levels & none (default) \\
$N$ & 35, 70 & 35, 50, 70 \\
\bottomrule
\end{tabular}
\end{table}

\textbf{Recommendation:} Single main simulation with key parameters,
separate one-at-a-time sensitivity analyses.


%% ============================================================
\section{Implementation Strategy}
%% ============================================================

\textbf{Phase 1: Minimal Viable Simulation.}
Single script \texttt{simulation.R} (${\sim}200$ lines); random
intercept model for data generation; ANCOVA on phase means for analysis;
four designs with simplified specification; no carryover, no dropout;
fixed BR\_rate, ER\_rate; vary biomarker\_mod and $N$.

\textbf{Phase 2: Add Complexity as Needed.}
Add carryover (linear decay); add dropout (MCAR); sensitivity on
BR\_rate; compare ANCOVA vs.\ LMM results.

\textbf{Phase 3: Validation.}
Compare power estimates to pmsimstats2025; document where
simplifications matter vs.\ do not.


%% ============================================================
\section{Code Architecture}
%% ============================================================

\begin{Verbatim}[fontsize=\small]
01b-pmsimstats-simple/
  R/
    functions.R            # Shared utility functions
  analysis/
    scripts/
      simulation.R         # Main simulation (single file)
      sensitivity.R        # Parameter sensitivity (optional)
    output/
  docs/
    simplification-plan.md
    simplification-plan.tex
\end{Verbatim}

\textbf{Key design principles:}
(1)~Single file first.
(2)~Explicit over clever: write out loops rather than complex
abstractions.
(3)~Base R preferred; minimize dependencies.
(4)~No parallelization by default.
(5)~Inline parameters at top of script.


%% ============================================================
\section{Expected Outcomes}
%% ============================================================

\begin{table}[h]
\centering
\begin{tabular}{@{}lll@{}}
\toprule
Metric & pmsimstats2025 & pmsimstats-simple (target) \\
\midrule
Lines of code & ${\sim}700$ per script & ${\sim}200$ total \\
Dependencies & nlme, lme4, MASS, zoo, patchwork, future, furrr & tidyverse only \\
Runtime (20 iter) & ${\sim}160$ sec & ${\sim}10$--20 sec \\
Convergence issues & Occasional lme failures & None (closed-form) \\
Parameters & ${\sim}15$ varying & ${\sim}5$ varying \\
\bottomrule
\end{tabular}
\end{table}


%% ============================================================
\section{Literature Support for Simplifications}
%% ============================================================

\subsection{Key Reference: Hendrickson et al.\ (2020)}

The primary reference is Hendrickson et al.\ (2020), ``Optimizing
Aggregated N-Of-1 Trial Designs for Predictive Biomarker Validation,''
\textit{Front.\ Digit.\ Health} 2:13.

Key methodological points:
\begin{enumerate}[nosep]
\item \textbf{Three-factor response model:} BR, ER, TR, each modeled
  via Gompertz function with max, rate, displacement parameters.
\item \textbf{Analysis approach:} LME (\texttt{lme4}) with random
  intercept per subject. For designs with on/off drug:
  \texttt{response \textasciitilde{} on\_drug + biomarker +
  on\_drug:biomarker + (1|id)}. For open-label only:
  \texttt{response \textasciitilde{} time + biomarker +
  time:biomarker + (1|id)}.
\item \textbf{Carryover:} Exponential decay with half-life parameter.
  Critical finding: even short (0.1 week) carryover causes
  ``precipitous decline in power'' for N-of-1 and OL+BDC designs.
\item \textbf{Censoring/dropout:} Probability =
  $\beta_0 + \beta_1 \times (\text{change in symptoms})^2$. Biased
  dropout where responders more likely to drop out during placebo.
\item \textbf{Key result:} N-of-1 design has power between OL+BDC
  (lower) and traditional crossover (higher), but allows all
  participants to start on open-label treatment.
\end{enumerate}

\subsection{Analysis Method Comparisons}

\textbf{Chen \& Chen (2014)} compared paired $t$-test, mixed effects
model, mixed effects model of difference, and meta-analysis for N-of-1
trials (3--4 cycles, $N = 1$--30):

\begin{table}[h]
\centering
\begin{tabular}{@{}lllll@{}}
\toprule
Method & Type I Error & Power & Bias & MSE \\
\midrule
Paired $t$-test & Near nominal & Highest & Comparable & Smallest \\
Mixed effects & Similar & Lower & Smaller & Bigger \\
ME of difference & Far from nominal & Low & Large & Large \\
Meta-analysis & Far from nominal & Low & Large & Large \\
\bottomrule
\end{tabular}
\end{table}

Recommendation: ``Paired $t$-test was recommended to use for normally
distributed data of N-of-1 trials irrespective with or without
carryover effect.''

\textbf{Critique (Araujo et al.\ 2016):} Chen \& Chen models ``do not
include a treatment by patient interaction'' advocated in meta-analysis
literature.

\subsection{ANCOVA for RCTs}

\textbf{BMC Medical Research Methodology (2014):} ``ANCOVA should be the
analysis of choice, a priori, for RCTs with a single post-treatment
outcome measure previously measured at baseline; its superiority is
particularly marked when baseline imbalance is present.''

\subsection{Biomarker-Treatment Interaction Power}

\textbf{Trials (2018):} ``Tests of interactions are often lacking
statistical power.'' ``The probability of detecting a biomarker-treatment
interaction can be increased by including prognostic variables.''

\subsection{N-of-1 Power Analysis}

\textbf{Wang \& Schork (2019):} Analytical and simulation-based power
analysis for N-of-1 designs considering number of cycles, serial
correlation, washout periods, heteroscedasticity, and carryover.

\subsection{Implications for Simplification}

\begin{table}[h]
\centering\small
\begin{tabular}{@{}L{6.5cm}L{7cm}@{}}
\toprule
Literature Finding & Implication for pmsimstats-simple \\
\midrule
Paired $t$-test has highest power (Chen 2014) & ANCOVA on change scores is reasonable \\
ANCOVA superior for baseline-adjusted outcomes (BMC 2014) & Supports phase-mean approach \\
Carryover devastating to N-of-1 power (Hendrickson 2020) & Must include carryover sensitivity \\
LME may have lower power than simpler methods & Justifies avoiding complex mixed models \\
Interaction tests lack power (Trials 2018) & Need adequate sample sizes \\
\bottomrule
\end{tabular}
\end{table}


%% ============================================================
\section{Decision Log}
%% ============================================================

\begin{table}[h]
\centering\small
\begin{tabular}{@{}L{3.5cm}L{8cm}l@{}}
\toprule
Decision & Rationale & Date \\
\midrule
Use ANCOVA & Chen 2014 shows paired $t$-test $\approx$ ANCOVA has highest power; BMC 2014 supports ANCOVA & 2026-02-14 \\
Keep 4 designs & Core scientific comparison per Hendrickson 2020 & 2026-02-14 \\
Keep ER component & Hendrickson shows expectancy affects power & 2026-02-14 \\
Include carryover sensitivity & Hendrickson shows carryover ``precipitous'' effect & 2026-02-14 \\
No carryover default & Simplifies initial analysis; add as sensitivity & 2026-02-14 \\
Random intercept + iid & Chen 2014 shows simpler models perform well & 2026-02-14 \\
\bottomrule
\end{tabular}
\end{table}


%% ============================================================
\section{References}
%% ============================================================

\begin{enumerate}[nosep]
\item Hendrickson RC, Thomas RG, Schork NJ, Raskind MA (2020).
  Optimizing Aggregated N-Of-1 Trial Designs for Predictive Biomarker
  Validation. \textit{Front.\ Digit.\ Health} 2:13.
  \url{https://doi.org/10.3389/fdgth.2020.00013}

\item Chen X, Chen P (2014). A Comparison of Four Methods for the
  Analysis of N-of-1 Trials. \textit{PLoS ONE} 9(2): e87752.
  \url{https://doi.org/10.1371/journal.pone.0087752}

\item Araujo A, Julious S, Senn S (2016). Understanding Variation in
  Sets of N-of-1 Trials. \textit{PLoS ONE} 11(12): e0167167.
  \url{https://doi.org/10.1371/journal.pone.0167167}

\item Egbewale BE, Lewis M, Sim J (2014). Bias, precision and
  statistical power of analysis of covariance in the analysis of
  randomized trials with baseline imbalance. \textit{BMC Med Res
  Methodol} 14:49. \url{https://doi.org/10.1186/1471-2288-14-49}

\item Eng KH, Konings P (2018). A simulation study on estimating
  biomarker-treatment interaction effects in randomized trials with
  prognostic variables. \textit{Trials} 19:72.
  \url{https://doi.org/10.1186/s13063-018-2491-0}

\item Wang Y, Schork NJ (2019). Power and Design Issues in
  Crossover-Based N-Of-1 Clinical Trials with Fixed Data Collection
  Periods. \textit{Healthcare} 7(3):84.
  \url{https://doi.org/10.3390/healthcare7030084}

\item Kravitz R, Duan N, editors (2014). Design and Implementation of
  N-of-1 Trials: A User's Guide. AHRQ Publication No.\
  13(14)-EHC122-EF.
  \url{https://effectivehealthcare.ahrq.gov/products/n-1-trials/research-2014-5}

\item Cohen J (1983). The cost of dichotomization. \textit{Applied
  Psychological Measurement} 7:249--253.

\item Royston P, Altman DG, Sauerbrei W (2006). Dichotomizing
  continuous predictors in multiple regression: a bad idea.
  \textit{Statistics in Medicine} 25:127--141.
\end{enumerate}

\vfill
\noindent\rule{\textwidth}{0.4pt}
{\footnotesize
Rendered on 2026-03-25 at 18:49 PDT.\\
Source: \textasciitilde/prj/alz/10-pmsimstats-ng/pmsimstats-ng/docs/17-simplification-plan.tex}

\end{document}
